%%%%%%%%%%%%%%%%%%%%%%%%%%%%%%%%%%%%%%%%%%%%%%%%%%%%%%%%%%%%%%%%%%%%%%%%
\section{Sketch of the lower bound}\label{sec:lower}
%%%%%%%%%%%%%%%%%%%%%%%%%%%%%%%%%%%%%%%%%%%%%%%%%%%%%%%%%%%%%%%%%%%%%%%%

\begin{theorem}\label{thm:lower}
For every $Q\ge2$ and fixed $\varepsilon>0$,
$\liminf_{n\to\infty}\log\Prob(N\ge n)/\log n\ge-4/\gamma^2$.
\end{theorem}

With Theorem~\ref{thm:upper} and Proposition~\ref{prop:critupper} this proves
Theorem~\ref{thm:tail}. We explain the structure and the key estimates; the complete argument is in
the files listed in Appendix~\ref{app:repo}. First let $Q>2$ and write
\[
\lambda_+=Q+\sqrt{Q^2-4}=\frac4\gamma,\qquad\lambda_-=Q-\sqrt{Q^2-4}=\gamma .
\]

\subsection{The difficulty with selecting a high square}

The heuristic says: find a square of large height, then count its descendants. But a high square
counts only if it is \emph{reachable}, that is, all its ancestors are active. And if one selects
``the first high square the subdivision finds'', the selection depends on the field around it, so
its descendants no longer have the simple Gaussian law of a fixed square. The proof instead counts
\emph{deterministic} squares, fixed in advance rather than found by the field, with three properties, each with a clean conditional probability
given the square's own average, and controls the count by a second moment.

\subsection{The candidates}

Fix a depth $j$. For each of the $4^j$ squares $S$ of depth $j$ consider the \emph{high band}
\[
B_S=\bigl\{h_S\in[\lambda_+j\log2,\;\lambda_+j\log2+1]\bigr\}.
\]
On $B_S$ the height of $S$ is about $y_j=\sqrt{Q^2-4}\;j\log2-Q\log\varepsilon$. Since
$\operatorname{Var}(h_S)=j\log2+O(1)$, the expected number $Z_j$ of squares in the high band is
\[
\E[Z_j]=2^{\,j(2-\lambda_+^2/2)+O(\log j)}=e^{-(\lambda_++o(1))y_j},
\]
where the last step uses $\lambda_+^2/2-2=\lambda_+(\lambda_+-Q)$, which is the statement
$\psi(\lambda_+)=0$. This is the cost $e^{-\lambda_+y}$ of the heuristic.

\emph{Reachability.} Conditional on $B_S$, the average over the ancestor of $S$ at depth $l$ is
normal with mean $\lambda_+l\log2+O(1)$ and variance at most $l\log2+O(1)$, uniformly. The
ancestor is active when this exceeds $Ql\log2+Q\log\varepsilon$, and $\lambda_+>Q$, so failure at
level $l$ has probability decaying exponentially in $l$. A compactness argument shows that all ancestors are
active with conditional probability at least some $c_0>0$ independent of $S$ and $j$. It uses a
\emph{Cameron--Martin shift}: adding to the field a smooth deterministic function, here one that
raises every relevant average by the same amount, changes the law of the field only by an explicit
density.

\emph{The ancestral ceiling.} This property keeps the ancestors of $S$ from being much higher than
needed. Fix a small $\varpi>0$ and let $\mathcal C_S$ be the event that the average
over every ancestor at depth $l\le j$ is at most $\lambda_+l\log2+\varpi j\log2$. Given $B_S$ the
ancestors follow their conditional means closely, so $\mathcal C_S$ fails with conditional probability
$o(1)$.

\emph{Many descendants.} $S$ has at least $e^{(\gamma-o(1))y_j}$ descendants all of whose ancestors
from $S$ down are active; given $B_S$ this has conditional probability tending to one
(Lemma~\ref{lem:desc}).

\subsection{The second moment and the role of the ceiling}

Let $T_j$ count the squares with $B_S\cap \mathcal C_S$. Sum over pairs $S,S'$ of depth-$j$ squares at distance
about $2^{-j_0}$; there are about $4^{2j-j_0}$ of them. Let $U$ be the ancestor of $S$ at depth $j_0$.
The averages over $S$, $S'$ and $U$ have covariances $j_0\log2+O(1)$ with each other, so
$h_S+h_{S'}-h_U$ has variance $(2j-j_0)\log2+O(1)$. If both squares are in the high band and the
ceiling holds at $U$, then $h_S+h_{S'}-h_U\ge\lambda_+(2j-j_0)\log2-\varpi j\log2$, and a
one-dimensional Gaussian bound gives
\[
\Prob(B_S\cap \mathcal C_S\cap B_{S'})\le C\,2^{-\lambda_+^2(j-j_0/2)+\lambda_+\varpi j}.
\]
Summing over $j_0$, each separation contributes at most $2^{\lambda_+\varpi j}\E[Z_j]$ up to constants,
because $2-\lambda_+^2/2<0$. Hence
\[
\E[T_j^2]\le\mathrm{poly}(j)\;2^{\lambda_+\varpi j}\;\E[Z_j].
\]
Without the ceiling, this second-moment bound holds only when $\lambda_+^2\ge8$, that is $c\le-2$.
For $c>-2$ the untruncated second moment is inflated by pairs of high squares whose common
ancestor at an intermediate depth is unusually high; the ceiling excludes such pairs, and it costs
only a conditional probability $o(1)$.

\subsection{The descendant lemma}

\begin{lemma}[descendants of a deterministic square]\label{lem:desc}
Let $Q\ge2$, $\lambda>Q$, and $0<\theta<\min(Q,2)$. For a deterministic square $S$ of depth $j$,
condition on $h_S\in[\lambda j\log2,\lambda j\log2+1]$, and let $y=(\lambda-Q)j\log2-Q\log\varepsilon$
be the resulting height. Then, uniformly in $S$ and in the conditioned value, with conditional
probability tending to one as $j\to\infty$, $S$ has at least $e^{(R_\theta-o(1))y}$ descendants
all of whose ancestors from $S$ downward are active, where
\[
R_\theta=\frac{2-\theta^2/2}{Q-\theta}.
\]
For $Q>2$ and $\theta=\gamma$, $R_\gamma=\gamma$.
\end{lemma}

\begin{proof}[Idea of the proof]
Let $\Phi_U=h_U-\E[h_U\mid h_S]$ for descendants $U$ of $S$; after rescaling $S$ to the unit square
these form a log-correlated family with bounded errors, uniformly in $S$. At a relative depth $j'$
slightly less than $y/((Q-\theta)\log2)$, give each descendant the weight
\[
w_U=4^{-j'}\exp\bigl(\theta \Phi_U-\theta^2\operatorname{Var}(\Phi_U)/2\bigr),
\]
and let $\mathcal M$ be the total weight. Three facts finish the argument.
(i) $\mathcal M$ is not small: all its negative moments are bounded uniformly, by Kahane's inequality with
the convex function $u\mapsto u^{-p}$ and the finiteness of negative moments of subcritical chaos
on a ball~\cite[Theorem~2.12]{RV}. So $\mathcal M\ge e^{-o(y)}$ with probability tending to one.
(ii) Bad descendants carry little weight: under the weighting, averages along the path to $U$
grow at slope $\theta$ while the threshold grows at slope $Q$, so by the choice of $j'$ the height
stays positive with room to spare, and descendants with an inactive ancestor or an unusually large
final value carry weight exponentially small in $y$.
(iii) Each good descendant has weight at most $2^{-j'(2-\theta^2/2)+o(y)}$. Dividing the good weight
by this maximum gives at least $e^{(R_\theta-o(1))y}$ good descendants. For $\theta=\gamma$, the
identity $2-\gamma^2/2=\gamma(Q-\gamma)$ gives $R_\gamma=\gamma$.
\end{proof}

Because the conditional probability tends to one, the lemma attaches to the candidates without
any conditional independence: given $B_S$, the intersection of the three properties has
probability at least $c_0-o(1)$ by the union bound.

\subsection{Assembly for \texorpdfstring{$Q>2$}{Q>2}}

Let $G_j$ count the depth-$j$ squares that are in the high band, satisfy the ceiling, are
reachable and active, and have at least $e^{(\gamma-o(1))y_j}$ active descendants. Then
$\E[G_j]\ge(c_0-o(1))\E[Z_j]$ and $G_j\le T_j$, so the Paley--Zygmund inequality gives
\[
\Prob(G_j>0)\ge\frac{\E[G_j]^2}{\E[G_j^2]}\ge\frac{(c_0-o(1))^2\,\E[Z_j]^2}{\E[T_j^2]}
\ge\exp\Bigl(-\Bigl(\lambda_++\frac{\lambda_+\varpi}{\sqrt{Q^2-4}}+o(1)\Bigr)y_j\Bigr).
\]
On $\{G_j>0\}$ the subdivision has at least $e^{(\gamma-o(1))y_j}$ squares. Monotonicity in $n$
fills in the sizes between consecutive values of $j$, and letting $\varpi\to0$ gives
$\liminf\log\Prob(N\ge n)/\log n\ge-\lambda_+/\gamma=-4/\gamma^2$.

\subsection{The critical endpoint}

At $Q=2$ the root $\lambda_+=2$ equals $Q$, and $\gamma=2$ is where subcritical chaos estimates
fail. Fix instead a slope $\lambda>2$ and a chaos parameter $\theta<2$. The reachability, ceiling
and second-moment estimates use only $\lambda>\max(Q,2)$, and the descendant lemma gives
$R_\theta=(2-\theta^2/2)/(2-\theta)=(2+\theta)/2$. The same assembly yields
\[
\liminf_{n\to\infty}\frac{\log\Prob(N\ge n)}{\log n}\ge-\frac{\lambda^2/2-2}{(\lambda-2)R_\theta}
=-\frac{\lambda+2}{2+\theta}.
\]
Letting $\lambda\downarrow2$ and $\theta\uparrow2$ gives $-1$. No critical chaos measure is needed.
